\chapter*{Ringraziamenti}
\pagestyle{empty}
\thispagestyle{empty}

Desidero innanzitutto ringraziare la mia relatrice, la prof.ssa Damiana Lazzaro, per avermi seguito con enorme costanza e dedizione durante i corsi da lei tenuti e per tutta la durata del tirocinio e della stesura della tesi.
Il supporto dato è stato fondamentale per produrre questo elaborato, unito a un sostegno umano sempre presente.
È inoltre riuscita a trasmettermi la sua passione per la matematica, utilizzata come strumento per comprendere a fondo le cose, e non come mezzo di tortura verso i suoi studenti.

Un pensiero importante va anche ai miei correlatori, il prof. Luigi Guiducci e la dott.ssa Giulia Paggi, e al resto del gruppo dell'INFN di Bologna, il dott. Filippo Mei e la dott.ssa Licia Mozzina.
Mi avete guidato con grande pazienza passo dopo passo all'interno del software di SND con supporto costante, dedicandomi molto tempo sebbene scarso per voi in primis.
Avete inoltre fatto di tutto per regalarmi l'esperienza indimenticabile vissuta al CERN tutti insieme: credo non ci siano parole per descrivere quanto accaduto in quei fantastici giorni, tra colpi di uncinetto nella sala conferenze e qualche birra ai barbecue.
Per tutti quelli che non erano presenti, giuro che ci sono stati anche dei momenti seri.

Un grazie di cuore ai miei genitori e ai miei fratelli, per tutto ciò che avete fatto per me in questi anni: senza il vostro supporto e incoraggiamento non sarei qui oggi.
A mia mamma Erica per comprendere le mie difficoltà e a mio padre Andrea per gli sproni continui.
A mio fratello Matteo, per essere anche il mio migliore amico, e a mia sorella Matilde, per voler stare con me fino all'ultimo minuto prima di tornare a Cesena.
Sappiate che questo traguardo è anche vostro.

\needspace{3\baselineskip}
Grazie a te, Ale, tu che mi hai supportato soprattutto in questo ultimo duro anno, tra un piantino di sfogo e un episodio di Stranger Things.
A furia di parlare di università, ormai conosci i nomi dei miei professori meglio di me; hai sopportato in prima persona ogni mio alto e basso di questo percorso, consolandomi quando ne avevo bisogno.
Ti ringrazio anche per il supporto che mi dai persino nelle piccole cose; forse gli invitati alla festa devono ringraziarti se siamo riusciti ad organizzare qualcosa oggi.

Come potrei non citarti, Gianluca: abbiamo iniziato il nostro percorso insieme nel 2018, quando ancora non sapevamo cosa fosse un'equazione (non posso dire un ciclo for, visto che tu eri già avanti).
Benché le nostre strade si siano separate, ci siamo comunque ritrovati ad aggiornarci tutti i giorni sui corsi che stavamo affrontando.
Con Emily, ci siamo tuffati a capofitto nel mondo dell'università, passando intere giornate immersi tra gli orientamenti e gli esercizi dei test d'ingresso.
Si potrebbe quasi dire che le mie scelte sono anche un po' colpa vostra (o merito, a seconda dei punti di vista).
Sappiate che mi dispiace per tutte le volte che, assieme ad Ale, Vincenzo, Viola, Giuseppe e Tommaso, vi ho tartassati di messaggi sul mio percorso universitario; credo che a questo punto sareste in grado di discutere la tesi al mio posto.

A Mattia, che è ormai il mio compagno di viaggio da quasi venti anni.
E lo intendo letteralmente, dai viaggi in bicicletta a quelli sull'Intercity 614 per ritornare alla realtà di tutti i giorni dopo il weekend passato a casa.
Se in questi anni mi sono trasformato, lo devo anche al tuo incoraggiamento, insieme a quello di David, Sheron e Martina.

I miei due più stretti unifriend, che a farlo apposta non ci sarei riuscito, sono Nicolas e Nicholas, aka Nico e Svitol.
Abbiamo affrontato fianco a fianco qualsiasi difficoltà, partendo dal primo esercizio di analisi fino ad arrivare all'ultimo progetto di gruppo; e di progetti ne abbiamo dovuti fare eh.
Senza di voi (e le piadinate fatte alla piadineria Irma) la laurea non avrebbe avuto lo stesso senso.

A Giulio, Lorenzo e Matteo -- ordine rigorosamente alfabetico per non far torto a nessuno --, con cui in questo anno da fuorisede abbiamo condiviso di tutto: piatti per la cena, sedie di casa orripilanti, notti passate a ripassare prima degli esami e chi più ne ha più ne metta.
Avete reso la lontananza da casa un peso molto più gestibile di quello che sarebbe stato senza di voi.

Un bacio va a voi nonni -- Duilio, Assunta, Eliseo e Lilliana -- per avermi sempre voluto bene, nutrito (forse un po' troppo) e sostenuto da quando ancora non sapevo camminare.

Anche voi, zii e cugini, ad ogni esame mi avete sempre chiesto come fosse andata e, nonostante la lontananza o gli impegni personali, siete sempre riusciti a trovare un modo per scambiare due chiacchiere.
Capite anche voi però che elencarvi tutti richiederebbe una tesi a parte!

Grazie infine a tutti voi che, sebbene non nominati nelle righe precedenti, mi siete stati accanto in tutti questi anni e state ora leggendo queste parole, o ascoltando se sarò in grado di recitarvele, ma direi che arrivati a questo punto la risposta si saprà già.

Vorrei dedicare un ultimo pensiero a me stesso, per aver sempre tenuto duro anche quando le cose hanno iniziato a farsi difficili; ma alla fine, mio padre mi ha sempre detto che quando il gioco si fa duro, i duri iniziano a giocare.

Finalmente lo possiamo dire, anche a chi non ci contava:\\
\#ciVediamoAllaLaurea.